\documentclass[10pt,a4paper]{amsart}
\usepackage[margin=19mm]{geometry}
\usepackage{amsmath,amssymb,mathtools}
\usepackage[scr=rsfs,cal=euler]{mathalfa}
\usepackage{xfrac}
\usepackage[unicode,hidelinks,bookmarks=false]{hyperref}
\newcommand{\ie}{i.e.\ }
\newcommand{\cf}{cf.\ }
\newcommand{\resp}{respectively\ }
\newcommand{\Subsection}{\subsection}
\providecommand{\nobreakdash}{\nobreak\hbox{-}}
\setlength{\emergencystretch}{2em}
\multlinegap0pt
\usepackage{amsrefs}
\usepackage{booktabs}
\usepackage{mathtools}
\usepackage{amscd}
\usepackage{enumitem}
\usepackage{tikz}
\usepackage{tcolorbox}
\usepackage{cleveref}
\usepackage{latexsym}
\usepackage{mathrsfs}
\usepackage{psfrag}
\usepackage{epsfig}
\usepackage{epsfig}
\usepackage{xy}
\setlength\parindent{0pt}
\newcommand\Tstrut{\rule{0pt}{3ex}}         % = `top' strut
\newcommand\Bstrut{\rule[-0.9ex]{0pt}{0pt}}   % = `bottom' strut
\newcommand{\mcC}{\mathcal{C}}
\newcommand{\mcP}{\mathcal{P}}
\newtheorem{thm}{Theorem}[section]
\newtheorem{prop}[thm]{Proposition}
\newtheorem{defn}[thm]{Definition}
\newtheorem{lem}[thm]{Lemma}
\newtheorem{cor}[thm]{Corollary}
\newtheorem*{conj}{Conjecture}
\newtheorem*{mainthm}{Main Theorem}               % Or Theorem A, etc.
\newtheorem{example}[thm]{Example}
\newtheorem*{remark}{Remark} 
\newtheorem*{eg}{Example}
\newtheorem*{rem}{Remark}                         % non-bold heading
\newtheorem{case}{Case}
\newtheorem*{note}{Note}
\numberwithin{equation}{section}
\renewcommand{\baselinestretch}{1.09}
\newcommand{\R}{\mathbb{R}}                     % Reals
\newcommand{\Z}{\mathbb{Z}}                     % Integers
\newcommand{\C}{\mathbb{C}}                     % Complex numbers
\newcommand{\N}{\mathbb{N}}                     % Natural numbers
\newcommand{\mcD}{\mathcal{D}}                     % Natural numbers
\newcommand{\ms}[1]{\mathscr{#1}}               % Script
\newcommand{\too}{\longrightarrow}
\newcommand{\vect}[1]{\mathbf{#1}}              % Vectors
\newcommand{\noi}{\noindent} 
\newcommand{\veps}{\varepsilon}        
\newcommand{\del}{\partial}
\DeclareMathOperator{\rank}{rank}
\DeclareMathOperator{\codim}{codim}
\DeclareMathOperator{\vol}{vol}
\newcommand{\aperta}{\setlength{\itemsep}{-1pt}} % This shrinks enumerate 
\begin{document}
\title[Sobolev Spaces and Their Applications to Partial Differentia]{Sobolev Spaces and Their Applications to Partial Differential Equations}
\author{Shihan Kanungo ạte page 1 e̱gin abstract; Shihan Kanungo}
\date{}
\maketitle
\noindent\emph{Source and licence.} Sobolev Spaces and Their Applications to Partial Differential Equations. Version arXiv:2606.14994v1 (CC BY 4.0). This material is distributed under the Creative Commons Attribution 4.0 International licence, http://creativecommons.org/licenses/by/4.0/, as recorded in the arXiv OAI licence metadata for this exact version and shown on the abstract page https://arxiv.org/abs/2606.14994v1. Full rights notice: licenses/arxiv-2606.14994-CC-BY-4.0.txt.\par\smallskip
\noindent\textit{Editorial scope.} Counted unit: the original \texttt{thm} environment \texttt{-} at source lines 294--300 together with its complete proof at lines 301--313; the delivered document keeps source lines 286--314 so that every referenced label and every citation resolves inside it. Native editable source is the paper's own concatenated TeX; the AMS wrapper is editorial formatting only. No publisher class, upstream script or unrelated artwork is redistributed.
\section{The Sobolev Embedding Theorem}
We can think of the Sobolev space $W^m(U)$ as the $L^2$ functions that are ``$L^2$-differentiable'' up to order $m$. This does not coincide with the ordinary notion of smoothness: for example the function $|x|$ is ``$L^2$-differentiable'' up to order $1$, but is not $C^1$. Thus, it makes sense to ask the following question: what can we say about the smoothness of elements of the Sobolev space $W^m(U)$?

This question is answered by the Sobolev Embedding Theorem, and the answer turns out to be quite nice.

We let $C_0$ denote the space of continuous functions $f$ on $\R^n$ to vanish at $\infty$; i.e. the space $\{x: |f(x)|>\varepsilon\}$ is compact for any $\varepsilon>0$. We let
\[C_0^k = \{f\in C_0: \partial^\alpha f \in C_0 \quad\forall |\alpha|<k\},\]
which means that $f$ has all (classical) derivatives up to order $k$ continuous and vanishing at $\infty$.

\begin{thm}[Sobolev Embedding Theorem]
    Suppose $m>k+n/2$. Then,
    \begin{enumerate}
        \item[$(1)$] If $f\in H^m$, then $\widehat{\partial^\alpha f}\in L^1$ and $\Vert \widehat{\partial^\alpha f}\Vert_1 \le C \Vert f \Vert_{(m)}$ for all $|\alpha|\le k$, where $C$ is a constant depending only on $m-k$.
        \item[$(2)$] $H^m \subset C_0^k$ and the inclusion map is continuous.
    \end{enumerate}
\end{thm}

\begin{proof}
    For $f\in H^m$ and $|\alpha|\le k$, we have
    \begin{align*}
        (2\pi)^{-|\alpha|} \Vert \widehat{\partial^\alpha f}\Vert_1 &=(2\pi)^{-|\alpha|}\int |(\widehat{\partial^\alpha f})(\xi) |\, d\xi\\  &= \int |\xi^\alpha \widehat{f}(\xi)|\, d\xi \\
        &\le \int (1+|\xi|^2)^{k/2} |\widehat{f}(\xi)|\, d\xi.
    \end{align*}
    This is because $|\alpha| \le k$ implies $|\xi|^{2\alpha}$ is a term in the expanded form of $(1+|\xi|^2)^k$, so $|\xi^\alpha|\le (1+|\xi|^2)^{k/2}$. Now applying the Cauchy-Schwarz inequality,
    \begin{align*}
        \int (1+|\xi|^2)^{k/2} |\widehat{f}(\xi)| \, d\xi &\le \left( \int (1+|\xi|^2)^{m} |\widehat{f}(\xi)|^2\, d\xi\right)^{1/2}\left( \int (1+|\xi|^2)^{k-m}\, d\xi\right)^{1/2}\\
        &= \Vert f \Vert_{(m)}\left( \int (1+|\xi|^2)^{k-m}\, d\xi\right)^{1/2}.
    \end{align*}
    Since $m-k< - n/2$, the last integral above is finite. Combining everything together, we have proven claim (1). The \textit{Riemann-Lebesgue Lemma} says that the Fourier transform of a $L^1$ function belongs to $C_0$. Using the Fourier inversion theorem, it follows that $\partial^\alpha f\in C_0$. From here, claim (2) follows.  
\end{proof}


\end{document}
